\section{Prototype}
\label{sec:prototype}

\begin{figure}[t]
    \centering
    \begingroup
    \setlength{\fboxsep}{3pt}
    \setlength{\tabcolsep}{2pt}
    \resizebox{\linewidth}{!}{%
    \begin{tabular}{c c c c c c c}
        \fbox{\parbox[c][0.52cm][c]{1.45cm}{\centering\scriptsize XIAO\\microphone}}
        & $\rightarrow$ &
        \fbox{\parbox[c][0.52cm][c]{1.45cm}{\centering\scriptsize Audio\\window}}
        & $\rightarrow$ &
        \fbox{\parbox[c][0.52cm][c]{1.55cm}{\centering\scriptsize Log-mel\\features}}
        & $\rightarrow$ &
        \fbox{\parbox[c][0.52cm][c]{1.55cm}{\centering\scriptsize Integer\\inference}} \\
        \multicolumn{7}{c}{$\downarrow$} \\
        \fbox{\parbox[c][0.52cm][c]{1.45cm}{\centering\scriptsize Intent +\\confidence}}
        & $\rightarrow$ &
        \fbox{\parbox[c][0.52cm][c]{1.45cm}{\centering\scriptsize Timing\\fields}}
        & $\rightarrow$ &
        \fbox{\parbox[c][0.52cm][c]{1.55cm}{\centering\scriptsize Safety-state\\bridge}}
        & $\rightarrow$ &
        \fbox{\parbox[c][0.52cm][c]{1.55cm}{\centering\scriptsize Gated Tello\\SDK command}}
    \end{tabular}%
    }
    \endgroup
    \caption{Prototype data path. The board performs local audio capture, feature extraction, and integer inference, then exposes intent-event fields to a safety-state bridge that gates Tello SDK command dispatch.}
    \label{fig:prototype_pipeline}
\end{figure}


\subsection{Embedded Voice-Interaction Event Source}

The prototype turns the recognizer into an embedded voice-driven safety event source. The hardware baseline is the XIAO ESP32-S3 Sense running local microphone capture, log-mel feature extraction, int8 TFLM inference, and USB CDC result reporting. The board does not issue flight commands; it emits the event fields needed by the control bridge. Figure~\ref{fig:prototype_pipeline} summarizes the data path.

\subsection{Onboard Inference Path}

A host trigger over USB CDC causes the ESP32 to capture a $1$ s onboard audio window, compute the log-mel frontend, run the embedded student as an int8 TFLM model, and return predicted label, confidence, timing, and raw output. The full-integer artifact is $780416$ bytes, and the exported operator mix includes one \texttt{DEPTHWISE\_CONV\_2D} for the temporal convolution. These details matter because the prototype is constrained by TFLM operator support as much as by model accuracy.
% Evidence: full-integer artifact and op mix from weeklyresult/weekly_drone_2026w17/B_small_teacher_student/tflm_candidate_precheck.json.

\subsection{Voice Event Output}

The embedded path emits a voice event rather than a direct control command. The event contains predicted label, confidence, inference timing, total timing, and raw-output diagnostics. A $30$-trigger USB CDC stability run records success rate $1.0000$, drop rate $0.0000$, inference p50 $2094.0$ ms, total p50 $3075.0$ ms, and raw uniform count $0$. These measurements characterize runtime stability for the local inference loop. They do not establish labeled real-world semantic accuracy.
% Evidence: weeklyresult/weekly_drone_2026w17/realworld/esp32_bench/local_cdc_fast_v4_stability30_report.md.
% TODO(validation): Add labeled real-world semantic accuracy before using these runtime logs as recognition-performance evidence.

\subsection{Control-Bridge Interface}

The bridge consumes voice events and produces logged state transitions. At minimum, each bridge record should include predicted label, confidence, safety state, selected command, command send event, and acknowledgement or timeout. Movement events require an additional policy source before any directional action, while emergency events are eligible for hold or operator escalation. The dry-run dispatcher exercises this interface without sending drone UDP packets, which makes it useful for validating state transitions before no-propeller or live experiments.
% Evidence: dry-run dispatcher evidence is weeklyresult/weekly_drone_2026w18/realworld/tello_dryrun/20260504_local_cdc_replay_report.md.
% TODO(validation): Add no-propeller bridge logs before claiming live control behavior.

The prototype keeps the bridge host-mediated for this draft. That choice is conservative: it lets the board remain a voice-event source while the host records bridge decisions, command eligibility, and acknowledgement behavior. A later design can move more of the bridge onto the vehicle, but the paper does not need that migration to argue for a voice-driven UAV safety mechanism. What matters is that the interface between recognition and control is explicit and logged.

\subsection{Trigger Gate and Buffer}

The architecture can add a board-side trigger gate and circular buffer before full inference so the student recognizer does not run continuously. The gate is an efficiency and false-trigger-control component, not the safety claim itself. It must be validated against silence, background, rotor-like noise, and speech-trigger cases before it becomes part of the reported mechanism.
% TODO(validation): Existing gate/buffer work is preliminary and should stay out of visible claims until false-trigger behavior is controlled.
